% Novelty index of the thesis --- Fix_Notes item on new results,
% how they differ from or extend previous work, and why they matter.
% Build: latexmk -pdf novelty_index.tex
\documentclass[11pt,a4paper]{report}

\usepackage[margin=22mm]{geometry}
\usepackage[T1]{fontenc}
\usepackage[utf8]{inputenc}
\usepackage{lmodern}
\usepackage{microtype}
\usepackage{amsmath,amssymb,mathtools}
\usepackage{booktabs}
\usepackage{longtable}
\usepackage{array}
\usepackage{tabularx}
\usepackage{enumitem}
\usepackage{xcolor}
\usepackage{csquotes}
\usepackage{setspace}
\usepackage[hidelinks,pdfborder={0 0 0}]{hyperref}
\hypersetup{pageanchor=true}
\usepackage[nameinlink,noabbrev]{cleveref}

\setstretch{1.08}
\setlist[itemize]{leftmargin=1.4em,itemsep=0.25em,topsep=0.4em}
\setlist[enumerate]{leftmargin=1.6em,itemsep=0.35em,topsep=0.4em}
\setcounter{tocdepth}{1}
\setlength{\LTpre}{6pt}
\setlength{\LTpost}{10pt}
\setlength{\tabcolsep}{4pt}

\newcommand{\yp}{\textit{Yersinia pestis}}
\newcommand{\ft}{\textit{Francisella tularensis}}
\newcommand{\Ifix}{\widehat{I}}
\newcommand{\cls}[1]{\textsc{#1}}
\newcommand{\orig}{\cls{Original}}
\newcommand{\ext}{\cls{Extends}}
\newcommand{\alr}{\cls{Already}}
\newcommand{\nar}{\cls{Not a result}}

% Item-index longtable preamble: ID | name | class | nearest literature
\newcommand{\itemhead}{%
  \toprule
  ID & Short name & Class & Nearest literature \\
  \midrule
  \endfirsthead
  \toprule
  ID & Short name & Class & Nearest literature \\
  \midrule
  \endhead
  \midrule
  \multicolumn{4}{r}{\footnotesize\itshape continued} \\
  \endfoot
  \bottomrule
  \endlastfoot
}

\newcolumntype{L}[1]{>{\raggedright\arraybackslash}p{#1}}
\newcommand{\itemtable}{\begin{longtable}{L{1.55cm}L{4.35cm}L{2.35cm}L{6.55cm}}}

\title{\textbf{Novelty index of the thesis}\\[6pt]
\large Results that are new, that extend previous work,\\
and that are already in the literature}
\author{Adam Aldridge}
\date{21 August 2026}

\begin{document}
\hypersetup{pageanchor=false}
\maketitle
\hypersetup{pageanchor=true}

\begin{quote}
\small
\noindent\textit{Fix\_Notes (Novelty).}
In all the results, both presented and derived, you must emphasize those
which are new, how they differ from or extend previous work, and a
discussion of why they are important.
\end{quote}

\tableofcontents
\clearpage

% =====================================================================
\chapter{How to read this index}
\label{ch:how}

This document is a literature map of every extracted result in the thesis
(320 items across ten chapters). It is not a rewrite of the chapters. Use
it to decide what the thesis should cite as known, what it should claim as
original, and why that claim matters.

Ten extractors listed every theorem, closed form, model, comparison, and
argument that functions as a result. Literature agents searched the
mathematical object (not the slogan) against the chapter bibliographies
and the wider literature. Adversarial checkers then tried to downgrade
marks of novelty. One checker opened Karlin and Tavar\'e (1982) in full.

\paragraph{Classes.}
\begin{description}[style=nextline,leftmargin=2.2cm,labelwidth=2.0cm]
\item[\orig]
This formula, identity, or named construction was \emph{not found} in the
literature. Neighbours exist; the method or the question may be published.
This is the strongest claim made here.
\item[\ext]
A published result is the ancestor; the thesis changes the process, the
object, the hypotheses, or the biological reading.
\item[\alr]
Equivalent statement already published, even if re-derived here.
\item[\nar]
Recap, definition, figure, or framing.
\end{description}

\noindent
``Not found'' is not a uniqueness proof. When evidence was thin, the class
was moved toward \alr{} or \ext, not toward \orig. After adversarial
checking, \textbf{no result is without published ancestors.} Several
statements were still \textbf{not found as written}. Those, together with
the ancestor and the increment, are what the Fix\_Notes item wants
emphasized.

% =====================================================================
\chapter{Whole-thesis verdict}
\label{ch:verdict}

The analytical spine of the thesis is the linear birth--death process with
per-capita total catastrophe (killing into a cemetery \(H\)). That process,
its Riccati / linear-fractional generating function, the geometric
occupancy, the geometric killing position, and the identity ``burst size
equals the quasi-stationary law'' are \textbf{Karlin and Tavar\'e
(1982)}. The constant \(A(p)\) of Chapter~3 is \textbf{Kolmogorov
(1938)}. Constant-rate BMVR, \(R_0\) at matched lifetime yield, and
stochastic bursting-versus-budding at establishment are \textbf{Perelson /
Nowak / McLean} and \textbf{Pearson, Krapivsky and Perelson (2011)}.
Antibody flooding is \textbf{Komarova (2007)}. Diversity-dependent and
time-dependent speciation alternatives to Nee's push of the past are a
large existing literature.

What the thesis actually adds is a \textbf{closed-form intracellular
calculus} on top of Karlin--Tavar\'e, a \textbf{renewal embedding} of those
kernels into the within-host model, an \textbf{\(L\)-organised
comparison} of bursting and budding, a \textbf{two-type killing
deformation} of Antal--Krapivsky, two \textbf{finite-budget immune
caricatures}, and a cluster of \textbf{evolutionary / plague-strategy}
models that are applications or unfinished specifications. Those are the
results to emphasize. The 1982 / 1938 / 2011 / 2007 facts must be cited as
such, or the novelty claims will not survive viva.

The four \orig{} statements in the strict sense (no published equivalent
formula) are CH6-R20, CH6-R21, CH7-R06, and CH7-R23. Adversarial checkers
would label even those \ext, because Pearson+Karlin--Tavar\'e and
Antal--Krapivsky supply the \emph{method}. In the thesis, write them as
original identities obtained by specialising those methods, not as a new
theory of branching.

% ---------------------------------------------------------------------
\section{Original statements (not found as written)}
\label{sec:original}

These are the closest the map comes to ``does not currently show up in the
literature.'' Each has a named neighbour. Write them as original
\emph{statements}, with the neighbour cited in the same sentence.

\itemtable
\itemhead
CH6-R20 &
Flooding identity \(\Delta z_{\mathrm{ext}}=(L-1)(1/m-1)\) and criterion
\(L>1\Leftrightarrow 1/\delta<1/\lambda+1/\mu\) &
\orig &
Pearson 2011 (arbitrary burst law vs geometric drip; no \(L\));
Quirouette 2023 (law, not timing) \\
CH6-R21 &
At \(L=1\) the bursting and budding offspring \emph{laws} coincide
(not just extinction probabilities); variance order flips with \(L-1\) &
\orig &
Pearson geometric drip; Karlin--Tavar\'e geometric kill-size; Quirouette \\
CH7-R06 &
Exact finite-time no-catastrophe \(S(t)\) for triangular two-type BD with
type-specific linear catastrophe \(\delta_1 X+\delta_2 Y\), Wronskian mix
of a real-axis \({}_2F_1\) pair &
\orig &
Antal--Krapivsky 2011 (same Gauss reduction, terminal-count PGF, no
catastrophe). \(\delta=0\) in the survival IVP does \emph{not} recover AK \\
CH7-R23 &
Killed (defective) bivariate PGF by the same mix; AK is this object at
\(\delta_i=0\) &
\orig &
AK 2011 (25)--(29) \\
CH4-R17--R20, R22 &
Finite-time dump \(W_t\): \(V(t)\), \(\mathbb{E}[W_t^2]\),
\(\mathrm{Var}(W_t)\), current yield \(\mathbb{E}[X_t+W_t]\) as functions
of \(I(t)\) &
\ext &
KT 1982 killing position is the \emph{eventual} geometric; not a
time-dependent dump coordinate \\
CH4-R23 &
Every load moment is a polynomial \(Q_p(I)\) of stated degree, roots
\(a,b\), leading coefficient \((-1)^p p!(\lambda/\delta)^p\) &
\ext &
Equivalent moments sit in the KT PGF; the named \(I\)-economy was not
found \\
CH5-R22 &
Telescoping / size-bias criterion for ``burst law = QSD'' &
\ext &
KT 1982 already has the identity; this is a new \emph{proof and
characterisation} \\
CH5-R26--R27 &
\(\mu=0\): \(\tau\mid\mathcal{K}=n\) harmonic / Gumbel limit &
\ext &
Classical Yule maxima; not written as a lytic-cell assay law \\
CH5-R31--R34 &
Multi-founder kernels \(g_k\); \(\mu=0\) yield \(k+\lambda/\delta\)
(sub-extensive) &
\ext &
Branching property is classical; this yield identity and the free-sum
trap were not found \\
CH5-R35--R42 &
Chained immediate-transfer BDC: negative-binomial rupture sizes, interval
Laplace transforms &
\ext &
Textbook sums of geometrics; the laboratory was not found under this
name \\
CH8-R03 &
Replicator--suppressor CTMC (birth vs unreplenished one-to-one
annihilation) &
\ext &
Pilyugin--Antia 2000 (finite killers that \emph{recycle}); Lanchester
linear law \\
CH8-R11--R12 &
\(\mathcal{R}(\vartheta)=V_\infty a^{-\vartheta}\), threshold
\(\vartheta^\ast=\log V_\infty/\log a\) &
\ext &
Geometric stop-loss is actuarial folklore; not found as a BDC + integer
immune budget \\
CH8-R17 &
Under threshold removal the ChPop flooding criterion \emph{reverses}
(\(L<1\) vs \(L>1\)) &
\ext &
Yuan--Allen 2011 reverse budding vs bursting under a mass-action immune
term (different mechanism); Komarova flooding is unconditional \\
\end{longtable}

% ---------------------------------------------------------------------
\section{Real extensions (ancestor plus increment)}
\label{sec:extensions}

\itemtable
\itemhead
CH3-R14, R17 &
\(A(p)=2\psi_{2p}(1/2)\); hypertranscendence of the logistic Koenigs
function as a corollary &
\ext &
Kolmogorov 1938 \(A\); Koenigs 1884; Becker--Bergweiler 1995 \\
CH3-R23, R24, R26 &
Binary product remainder, hybrid scheme, negative elementary-form search &
\ext &
Pollak 1969, Seneta 1968, Pakes 2026; Imomov's candidate fails on discrete
binary \\
CH5-R23--R24, R29 &
Three conditional burst means; late-burst size-bias &
\ext &
Geometric QSD (KT). Assay: late bursts are larger \\
CH5-R45--R48 &
Single-cell BDC vs budding table; \ft{} / anthrax numerical rows &
\ext &
Pearson 2011 population comparison with \emph{fixed} \(N\) \\
CH6-R04--R06 &
Nelson's age-structured system forced by closed-form BDC kernels (no free
\(P(a)\)) &
\ext &
Nelson et al.\ 2004; McKendrick--von Foerster \\
CH6-R09--R10 &
Exponential-phase map \((p_{\mathrm{eff}}(r), d_{I,\mathrm{eff}}(r))\)
with BDC endpoints &
\ext &
Euler--Lotka; Diekmann; Wallinga--Lipsitch. Fitted \(p\) is a phase
property \\
CH6-R16 &
Concrete 3-to-2 identifiability obstruction &
\ext &
Nelson 2004; Miao et al.\ 2011 \\
CH6-R22--R24 &
Growth--extinction tradeoff / Pareto in \(\delta\) (numerical) &
\ext &
Pearson slower burst establishment; Komarova budding spreads faster \\
CH7-R18 &
\(\lambda_2=0\) ordinary Bessel mix &
\ext &
AK modified-Bessel bi-critical (different ODE) \\
CH8-R05--R06 &
Certainty diagonal \(\pi_{j,j}\); rate-independent \(i>j\Rightarrow\)
escape &
\ext &
Pilyugin's \(P^\ast\) depends on rates because killers recycle \\
CH8-R18 &
Distributional flooding vs batch-count flooding &
\ext &
Komarova is the batch-count argument \\
CH9-R03 &
Two-layer unused-\emph{individual}-capacity clock on a Yule process &
\ext &
Kendall 1948 TD-Yule; Sepkoski; Rabosky/Etienne DDD on \emph{species}
count \\
CH10-R06, R08 &
BDC macrophage burst \(\to\) slightly supercritical BD \(\to\) host-to-host
GW &
\ext &
Antia 2003; Pearson; Gilchrist nested WH--BH; Kendall \\
CH2-R24 &
Gated release with opening intensity \(\eta A_t\) &
\ext &
Davis PDMP; Anick--Mitra--Sondhi on/off \\
\end{longtable}

% ---------------------------------------------------------------------
\section{What not to claim as new}
\label{sec:traps}

These are the statements the chapters (or Chapter~1's roadmap) are most
likely to over-claim. Cite them as previous work and move the emphasis to
the increment above.

\itemtable
\itemhead
--- &
The BDC process with cemetery \(H\) &
\alr &
Karlin and Tavar\'e 1982. Brockwell--Gani--Resnick 1982 is the
\emph{more general} neighbour (immigration, random catastrophe sizes) \\
--- &
Closed forms \(I(t)\), \(D(t)\), \(\Ifix\), \(J(t)\), \(K(t)\), roots
\(a,b\) &
\alr &
Karlin--Tavar\'e 1982 linear-fractional PGF \\
--- &
Occupancy geometric at every \(t\); QSD geometric; burst size = QSD &
\alr &
KT 1982 (2.8)--(2.9); Zhang--Zhu 2011 uniqueness. Chapter~5 headlines are
1982; the telescoping \emph{proof} is the increment \\
--- &
Mean productive lifetime \(\lambda^{-1}\log(a/(a-1))\) &
\alr &
KT Green function \\
--- &
\(A(p)=\lim S_n/(2p)^n\); \(S_n\sim 2/n\) at criticality; Yaglom mean
\(1/A\) &
\alr &
Kolmogorov 1938; Yaglom 1947. Kolmogorov is never cited in Chapter~3 \\
--- &
Constant-\(p\) BMVR &
\alr &
Perelson 1996; Nowak--Bangham 1996; McLean 1993 \\
--- &
\(R_0\) unchanged by bursting vs dripping at matched lifetime yield &
\alr &
Pearson 2011, opening paragraph \\
--- &
Age-structured replacement of \(pI\) &
\alr &
Nelson et al.\ 2004 (as a \emph{system}). The increment is BDC kernels,
not the renewal skeleton \\
--- &
Comparing bursting and budding &
\alr &
Pearson 2011; Komarova 2007; Yuan--Allen 2011 \\
--- &
Flooding as a biological idea &
\alr &
Komarova 2007 (replace the flooding placeholder cite) \\
--- &
Finite immune budget / inoculum escapes killing &
\alr &
Pilyugin--Antia 2000; Antia--Koella 1994; inoculum effect; complement
consumption (as a \emph{mechanism}) \\
--- &
``Nobody has published a logistic speciation model'' &
False &
Sepkoski; Rabosky DDL; Etienne--Haegeman DDD; Kendall TD-Yule \\
--- &
A no-death Yule tests Nee's push of the past &
Cannot &
POTPa needs extinction and conditioning on survival \\
--- &
Neutrophil-fraction threshold as the evolved \yp{} switch &
\alr &
Ke, Chen and Yang 2013: switch is temperature / pH / YCV, not \(p_N\).
The displayed formula also has a sign error \\
--- &
Chapter~1: second moment \(K\) ``beyond the current literature'' &
Not sustained &
Equivalent to \(\partial_{zz}G(1,t)\) from KT 1982 \\
--- &
Chapter~2 hypergeometric integral identity &
\alr &
Gradshteyn--Ryzhik 3.194.1 / DLMF 15.6.1 \\
--- &
Chapter~2 ABD generating function as ``a result of this thesis'' &
\alr &
Kendall interior PGF integrated against exponential absorption \\
\end{longtable}

% ---------------------------------------------------------------------
\section{Counts (reconciled)}
\label{sec:counts}

\begin{center}
\small
\begin{tabular}{@{}lrrrrr@{}}
\toprule
Chapter & Items & Original & Extends & Already & Not a result \\
\midrule
1 Introduction & 26 & 0 & 1 & 14 & 11 \\
2 Mathematical introduction & 58 & 0 & 7 & 47 & 4 \\
3 Galton--Watson and \(A(p)\) & 28 & 0 & 5 & 22 & 1 \\
4 BDC core & 42 & 0 & 5 & 28 & 9 \\
5 BDC distribution & 49 & 0 & 26 & 22 & 1 \\
6 Renewal / population & 33 & 2 & 23 & 5 & 3 \\
7 Two-type BDC & 31 & 2 & 13 & 8 & 8 \\
8 Bursting and budding pathogenesis & 21 & 0 & 13 & 5 & 3 \\
9 Variability in evolution & 17 & 0 & 9 & 5 & 3 \\
10 Conclusion, \yp{} & 15 & 0 & 7 & 6 & 2 \\
\midrule
\textbf{Total} & \textbf{320} & \textbf{4} & \textbf{109} & \textbf{162} & \textbf{45} \\
\bottomrule
\end{tabular}
\end{center}

% ---------------------------------------------------------------------
\section{Citations the thesis is missing or under-using}
\label{sec:cites}

Add these where the corresponding result is discussed. They are the
difference between ``we derived \(I(t)\)'' and ``Karlin--Tavar\'e gave
\(I(t)\); we add \(W_t\) and the \(I\)-polynomials.''

\begin{itemize}
\item \textbf{Karlin and Tavar\'e (1982)}, \emph{Linear birth and death
  processes with killing}, J.\ Appl.\ Prob.\ 19:477--487 --- the missing
  citation for Chapters~4 and~5.
\item \textbf{Zhang and Zhu (2011)} --- uniqueness of the geometric QSD
  under dual absorption.
\item \textbf{Kolmogorov (1938)} --- Chapter~3's \(A(p)\) and
  \(S_n\sim 2/n\).
\item \textbf{Yaglom (1947)}; Harris (1963); Athreya--Ney (1972).
\item \textbf{Pearson, Krapivsky and Perelson (2011)} --- already cited;
  must be credited with \(R_0\) invariance and the establishment
  comparison, not only ``the question.''
\item \textbf{Yuan and Allen (2011)} --- must be cited in Chapter~8.
  Published ranking-switch of budding vs bursting once an immune term is
  present.
\item \textbf{Komarova (2007)} --- replace the flooding placeholder.
\item \textbf{Nelson, Gilchrist, Coombs, Hyman and Perelson (2004)} ---
  the renewal/age-structured system Chapter~6 specialises.
\item \textbf{Miao, Xia, Perelson and Wu (2011)} --- structural
  non-identifiability.
\item \textbf{Pilyugin and Antia (2000)}; \textbf{Antia and Koella
  (1994)} --- finite nonspecific killers (they recycle; that is why their
  threshold depends on rates).
\item \textbf{Nee, May and Harvey (1994)}; \textbf{Rabosky and Lovette
  (2008)}; \textbf{Etienne and Haegeman (2012)}; \textbf{Pannetier et
  al.\ (2021)} --- Chapter~9 cannot claim uniqueness of a variable-rate
  alternative to POTPa.
\item \textbf{Ke, Chen and Yang (2013)}; \textbf{Cavanaugh and Randall
  (1959)} --- the \yp{} switch is temperature/F1/Yop, not a neutrophil
  fraction.
\item \textbf{Pakes (2026)} on the Kolmogorov constant --- no known
  explicit formula; Imomov's candidate fails on discrete binary.
\end{itemize}

% ---------------------------------------------------------------------
\section{Suggested sentence shapes}
\label{sec:sentences}

Do not write ``we obtain the no-rupture probability \(I(t)\).'' Write:

\begin{quote}
Karlin and Tavar\'e (1982) already give the linear-fractional generating
function of the linear killing process, and with it \(I(t)\), the
geometric occupancy, and the geometric killing position. What is
developed here is the finite-time dump \(W_t\) and the organisation of
every load moment as a polynomial in \(I\). Those are the kernels that
later chapters feed into a within-host renewal system.
\end{quote}

Do not write ``we compare bursting and budding.'' Write:

\begin{quote}
Pearson et al.\ (2011) showed that at matched lifetime yield the
mean-field kinetics coincide while stochastic extinction need not.
Substituting the BDC burst law yields the identity
\(\Delta z_{\mathrm{ext}}=(L-1)(1/m-1)\), which Pearson's fixed-\(N\)
comparison cannot flip. Under a finite immune toll the same \(L\)
criterion reverses (cf.\ Yuan and Allen 2011 for a different immune
mechanism).
\end{quote}

Do not write ``we solve the two-type BDC.'' Write:

\begin{quote}
Antal and Krapivsky (2011) solved the triangular two-type birth--death
process for terminal counts. Here the same Gauss reduction is applied to
the no-catastrophe probabilities under type-specific linear killing. The
objects differ: setting \(\delta_i=0\) in \(S,G\) does not recover their
generating functions; it recovers them from the killed PGF.
\end{quote}

% =====================================================================
\chapter{Chapter 1 --- Introduction}
\label{ch:1}

\paragraph{Source.} \texttt{CH1/document MAIN/sections/}.
26 items. No numbered theorems.

\paragraph{Scope.}
Frames early \yp{} infection as a stochastic small-count problem; reviews
branching and introduces BDC; sketches biphasic plague pathogenesis,
BMVR, bursting vs budding (including Komarova flooding); states four
research questions; roadmaps later chapters.

\paragraph{Verdict.}
This chapter \emph{presents} results that later chapters derive. Almost
nothing here is original. Its job, for novelty, is to name the ancestors
correctly so later claims of extension have a place to sit.

\paragraph{Emphasize (weakly).}
CH1-R10 is not new mathematics (Kendall/Haldane/Antia: establishment near
criticality is sensitive to founder number). As a \emph{reading} of the
macrophage phase it is the biological motivation for the BDC cell. Cite
Antia et al.\ (2003) and Cavanaugh/Ke; do not call the branching identity
a result of the thesis. The four research questions (R18--R21) are the
right organisers; each already has a literature.

\paragraph{Do not claim.}
Agricultural peak-shift (R01) as a result (the \(\sim 6000\)-year split
is published; the saddle-crossing gloss is an uncited hypothesis). BDC as
a process introduced here (R07). Flooding (R14; Komarova 2007). BMVR
(R12). Biphasic pathogenesis (R08--R09; Cavanaugh 1959; Pujol--Bliska;
Ke 2013). The roadmap's remark that the second moment of BDC ``extends
beyond the current literature'' (R23) is not sustained: \(K\) is in the
Karlin--Tavar\'e PGF.

\itemtable
\itemhead
R01 & Agricultural emergence / peak-shift & \ext\ (hypothesis) &
Achtman 1999; Rasmussen 2015; Rascovan 2019; Newman--Cohen--Kipnis 1985 \\
R02 & Early-infection stochasticity & \alr &
Standard small-count branching / epidemic establishment \\
R03 & Brownian encounters & \nar & --- \\
R04 & GW / Yule / BD recap & \alr & Galton--Watson; Yule; Kendall \\
R05 & Critical / super / sub & \alr & Athreya--Ney \\
R06 & Near-critical sensitivity & \alr & Kendall \((\mu/\lambda)^n\); Haldane \\
R07 & BDC introduced & \alr &
Brockwell--Gani--Resnick 1982; Karlin--Tavar\'e 1982 \\
R08 & Biphasic YP pathogenesis & \alr &
Cavanaugh and Randall 1959; Pujol and Bliska 2005; Ke et al.\ 2013 \\
R09 & F1 / Yop / Pla & \alr & Du 2002; Sebbane 2006; Ke 2013 \\
R10 & Macrophage as near-critical founder boost & \alr\ (as maths) &
Antia et al.\ 2003; Hataye 2019; Kendall \\
R11 & HIV as comparator & \alr & Perelson 1996; Hataye 2019 intermittency \\
R12 & BMVR & \alr & Nowak 1996; Perelson 1996; McLean 1993 \\
R13 & Bursting vs budding & \alr &
Pearson 2011; Komarova 2007; Yuan--Allen 2011 \\
R14 & Flooding advantage & \alr & Komarova 2007. Must cite here \\
R15 & Eclipse / multiscale aside & \alr &
Nelson 2004; Rong; Guedj; Ciupe 2024. Sets up Ch.~6 \\
R16 & Punctuated equilibria motivation & \nar & Eldredge--Gould \\
R17 & Push of the past & \alr & Nee et al.\ 1994 \\
R18 & RQ1: why biphasic YP & \nar & Question \\
R19 & RQ2: bursting advantage & \nar & Question \\
R20 & RQ3: BMVR intracellular update & \nar & Question \\
R21 & RQ4: variable-rate evolution & \nar & Question \\
R22 & Roadmap Ch.~2 & \nar & --- \\
R23 & Roadmap BDC ``new second moment'' & \alr\ (claim false) &
Karlin--Tavar\'e. Over-claim. Fix. \\
R24 & Roadmap two-type / chained & \nar & Points at Ch.~5/7 \\
R25 & Roadmap bursting--budding models & \nar & Points at Ch.~6/8 \\
R26 & Roadmap evolution / conclusion & \nar & Points at Ch.~9/10 \\
\end{longtable}

\paragraph{Why it matters.}
It is the only place the four questions sit together. For novelty, its
contribution is attribution: every later original statement needs the
ancestor named here so the increment is visible.

% =====================================================================
\chapter{Chapter 2 --- Mathematical introduction}
\label{ch:2}

\paragraph{Source.} \texttt{CH2/sections/} plus appendices A--D.
58 items.

\paragraph{Scope.}
Methods: discrete and continuous Markov chains, branching, linear
birth--death, BDC as Def.~2.4, method of characteristics, small-population
examples, absorption models, a hypergeometric integral, coefficient
extraction. The chapter flags two items as ``of this thesis'': the
absorption--birth--death generating function and the hypergeometric
identity. Neither is original.

\paragraph{Verdict.}
No original statements. This is a toolkit chapter. Do not let worked
examples (gated release, log-spectral mean, ABD) drift into ``we
introduce the process.''

\paragraph{Emphasize.}
Nothing as new mathematics. Optionally, as worked models:
gated release (R24), opening intensity \(\eta A_t\), sits in Davis's PDMP
class; published on/off gates are autonomous. Log-spectral /
logistic-speciation mean (R18, R20--R22) is Kendall time-inhomogeneous
Yule with an individual-capacity clock; a preview of Chapter~9, not DDD.

\paragraph{Do not claim.}
ABD generating function (R49) as a result of this thesis (Kendall
interior PGF integrated against exponential absorption, rewritten by
GR~3.194.1). Hypergeometric identity (R56) as original (GR~3.194.1 /
DLMF~15.6.1). BDC (R12) and \(W_t\) (R13) as a new process
(Karlin--Tavar\'e 1982).

\itemtable
\itemhead
R01 & Chapman--Kolmogorov & \alr & Norris; Karlin--Taylor; Allen \\
R02 & Simple random-walk law & \alr & Feller; Karlin--Taylor \\
R03 & Birth--death embedded walk & \alr & Standard embedding \\
R04 & Gambler's ruin & \alr & Feller \\
R05 & Galton--Watson process & \alr & Harris; Athreya--Ney \\
R06 & Binary GW survival iteration & \alr & Textbook \\
R07 & Competing clocks & \alr & Exponential races \\
R08 & Gillespie direct method & \alr & Gillespie 1977 \\
R09 & Poisson process law & \alr & Textbook \\
R10 & Yule process law & \alr & Yule; Kendall \\
R11 & Linear birth--death process & \alr & Kendall 1948 \\
R12 & Birth--death--catastrophe process & \alr &
Karlin--Tavar\'e 1982 linear killing. Chapter over-claims if read as new \\
R13 & Released count and burst size & \alr &
KT Lemma~1, killing position \(K=X(T_H)\) \\
R14 & Non-catastrophe factorisation & \alr & Branching property \\
R15 & Catastrophe is not mean-field & \alr & Jensen / convexity caution \\
R16 & Inhomogeneous Kolmogorov & \alr & Standard \\
R17 & Inhomogeneous linear BD mean and extinction & \alr & Kendall \\
R18 & Seasonal Lotka--Volterra process & \ext & Classical LV as illustration \\
R19 & Seasonal LV simulation & \nar & Figure \\
R20 & Logistic-speciation model & \ext & Kendall TD-Yule; preview of Ch.~9 \\
R21 & Log-spectral mean closed form & \ext &
Kendall \(\mathbb{E}=\exp\int\beta\); explicit limit
\((C/N_0)^{\sigma C/r}\) \\
R22 & Logistic BD variant / turnover & \ext & Same clock \\
R23 & ODE--CTMC couplings and PDMP generator & \alr & Davis 1993 \\
R24 & Gated release process \(\eta A_t\) & \ext &
Davis PDMP; Anick--Mitra--Sondhi autonomous on/off \\
R25 & Gated balance and non-closed mean & \alr & \\
R26 & Surging gated release & \alr & \\
R27 & Gated vs surging ensemble & \nar & Figure \\
R28 & Coupled-system simulation & \alr & \\
R29 & First-step principle & \alr & Feller \\
R30 & Branching extinction fixed point & \alr & Harris \\
R31 & Forward/backward master & \alr & Kolmogorov \\
R32 & Birth--death PGF PDE & \alr & Kendall \\
R33 & Kendall birth--death generating function & \alr & Kendall 1948 \\
R34 & Birth--death mean & \alr & \\
R35 & Birth--death variance & \alr & \\
R36 & Critical infinite mean absorption time & \alr & \\
R37 & Birth--death extinction probability & \alr & \((\mu/\lambda)^n\) \\
R38 & Subcritical survival amplitude & \alr & Continuous-time Kolmogorov \\
R39 & CT conditional mean and QS limit & \alr & Yaglom / Kendall \\
R40 & Continuous-time amplitude in \(p\) & \alr & \\
R41 & Discrete amplitude and two constants & \ext &
Comparison of two classical survival amplitudes, not a new constant \\
R42 & Method of characteristics & \alr & PDE method \\
R43 & Absorption--death generating function & \alr & Linear PGF \\
R44 & Early-generation survival near criticality & \alr & \\
R45 & \(k\)-founder survival and establishment & \alr & Branching property \\
R46 & Cohort longevity / criticality sensitivity & \nar & Reading \\
R47 & Absorption-only binomial PGF & \alr & Binomial \\
R48 & ABD model and Riccati characteristic & \ext &
Two-type CTBP degeneration (Kendall interior + absorption) \\
R49 & ABD generating function & \alr &
Kendall + exponential absorption + GR 3.194.1. Chapter over-claims \\
R50 & ABD generating-function checks & \nar & Internal \\
R51 & ABD convolution representation & \alr & Standard first-event
decomposition \\
R52 & \(\psi\) hypergeometric form & \alr & GR / DLMF \\
R53 & ABD domain of validity & \alr & Domain of R49 \\
R54 & ABD resonance & \alr & \\
R55 & ABD coefficient recovery & \alr & \\
R56 & Hypergeometric integral identity & \alr &
GR 3.194.1 / DLMF 15.6.1. Not a result of this thesis \\
R57 & Coefficient extraction (Maclaurin / Cauchy) & \alr & Standard \\
R58 & Multilinear trinomial probabilities & \alr & Multinomial theorem \\
\end{longtable}

\paragraph{Why it matters.}
It supplies the notation and the BDC definition later chapters use.
Novelty for the thesis does not live here. The danger is the opposite:
presenting standard material, or table integrals, as derived results.

% =====================================================================
\chapter{Chapter 3 --- Galton--Watson and \texorpdfstring{\(A(p)\)}{A(p)}}
\label{ch:3}

\paragraph{Source.} \texttt{CH3/chapterB.tex}. 28 items.

\paragraph{Scope.}
Binary Galton--Watson (\(p_2=p\), \(p_0=1-p\)). Survival recursion;
subcritical Yaglom law with mean \(1/A(p)\); existence of
\(A(p)=\lim S_n/(2p)^n\); identification with the Koenigs function of the
logistic map; Becker--Bergweiler hypertranscendence; critical
\(S_n\sim 2/n\); a hybrid scheme for computing \(A(p)\); a negative
search for an elementary formula.

\paragraph{Verdict.}
No original statements. Kolmogorov is never cited, and most of the
chapter is Kolmogorov (1938) / Yaglom (1947) / Harris. What to emphasize
is a dictionary and a computational increment, plus a useful negative
experiment (Imomov's formula fails on discrete binary).

\paragraph{Emphasize.}
\begin{enumerate}
\item \(A(p)=2\psi_{2p}(1/2)\) (R14). Koenigs 1884 linearisation of the
  logistic map, evaluated at \(1/2\), is Kolmogorov's constant. Not a new
  object: a dictionary. Matters because a closed form for \(A\) would
  require a closed form for \(\psi\).
\item Hypertranscendence of \(z\mapsto\psi_r(z)\) (R17). Corollary of
  Becker--Bergweiler 1995. Closes the ``formula for \(\psi\), hence for
  \(A\)'' route. Does \emph{not} prove that \(p\mapsto A(p)\) is
  non-elementary.
\item Truncation bound and hybrid scheme (R23--R24). Pollak 1969 /
  Seneta 1968 / Pakes 2026 already iterate \(\mu_n=m^n/S_n\). The binary
  product remainder and the \(2{-}4p\) switch are the increment.
\item Negative elementary-form search (R26). Pakes (2026): no known
  explicit expression. Imomov--Murtazaev's candidate holds for
  linear-fractional offspring and for continuous-time quadratic, and
  \emph{fails} for discrete binary (at \(p=0.3\): \(2.25\) vs \(2.825\)).
  The CT cousin \emph{is} elementary; binary \(A(p)\) is not known to be.
\end{enumerate}

\paragraph{Do not claim.}
Existence of \(A(p)\) (R07, R11) --- Kolmogorov 1938. Yaglom mean
\(1/A(p)\) (R08--R09). Critical \(S_n\sim 2/n\) (R19--R21). That the
chapter ``introduces'' \(A(p)\). \textbf{Fix:} cite Kolmogorov (1938) in
the first paragraph that defines \(A(p)\).

\itemtable
\itemhead
R01 & Binary GW model & \alr & Galton--Watson; Harris; Athreya--Ney \\
R02 & Unconditional mean \((2p)^n\); regimes & \alr & Textbook \\
R03 & Offspring PGF; extinction recursion & \alr & Harris \\
R04 & Survival recursion & \alr & \\
R05 & Ultimate survival / extinction certainty & \alr &
\(p\le 1/2\Rightarrow\eta=1\) \\
R06 & Finite-\(n\) conditional mean & \alr & \\
R07 & \(A(p)\) definition, \(S_n\sim A(p)\,(2p)^n\) & \alr &
\textbf{Kolmogorov 1938} \\
R08 & QS mean \(1/A(p)\) & \alr & Yaglom \\
R09 & Yaglom-limit existence & \alr & Yaglom 1947 \\
R10 & Limiting conditional variance & \alr & Harris / Athreya--Ney \\
R11 & Infinite-product existence, bounds & \alr & Kolmogorov; telescoping \\
R12 & Logistic rescaling of survival & \alr &
Logistic map \(x\mapsto 2px(1-x)\) is the binary GW PGF \\
R13 & Koenigs linearisation & \alr & Koenigs 1884 \\
R14 & \(A(p)=2\psi_{2p}(1/2)\) & \ext &
Dictionary Kolmogorov \(\leftrightarrow\) Koenigs \\
R15 & BB hypertranscendence of the inverse & \alr &
Becker--Bergweiler 1995 \\
R16 & Inverse preserves DA & \alr & Standard lemma (Fernandes 2021) \\
R17 & HT of the logistic Koenigs function & \ext & Corollary of R15--R16 \\
R18 & Elementary Koenigs at \(r=2,4\) & \alr & Classical \\
R19 & Critical \(S_n\sim 2/n\) & \alr & Kolmogorov 1938 \\
R20 & Infinite expected lifetime at criticality & \alr & \\
R21 & Critical power-law tails & \alr & \\
R22 & Near-critical \(A\sim 2\varepsilon\) & \alr & Kolmogorov amplitude \\
R23 & Truncation bound / stopping rule & \ext & Pollak / Seneta / Pakes \\
R24 & Product--asymptotic hybrid & \ext &
Same line; binary remainder is the increment \\
R25 & Complex-plane Koenigs geometry & \nar & Visual \\
R26 & Elementary-formula search (negative) & \ext &
Pakes 2026; Imomov fails on binary \\
R27 & Unconditional factorial / second moments & \alr & Harris \\
R28 & Near-critical moment rates & \alr & Yaglom + Kolmogorov amplitude \\
\end{longtable}

\paragraph{Why it matters.}
Binary GW is the discrete counterpart of linear birth--death, and
\(A(p)\) is the discrete cousin of the Karlin--Tavar\'e root \(b\). The
chapter's genuine service is to show why the discrete constant is harder
than the CT one, and to give a working numerical scheme.

% =====================================================================
\chapter{Chapter 4 --- BDC: definition and main results}
\label{ch:4}

\paragraph{Source.} \texttt{CH4\_REWRITE\_RESTORED/sections/}. 42 items.

\paragraph{Scope.}
Linear BDC as one lytic infected cell: per-capita rates
\(\lambda n,\mu n,\delta n\), cemetery \(H\), dump coordinate \(W_t\).
Riccati for \(I(t)\); roots \(a,b\); \(I,D,\Ifix\); load moments
\(J,K\); release moments of \(W_t\); burst-size first two moments;
\(I\)-polynomial organisation.

\paragraph{Verdict.}
This is the chapter most at risk of over-claiming. Karlin and Tavar\'e
(1982) already contain the process, the Riccati / linear-fractional PGF,
\(I(t)\), the geometric killing position with mean \(a/(a-1)\), and the
geometric QSD. Dictionary: their \((a,b,c)\) (death, birth, killing) are
this chapter's \((\mu,\lambda,\delta)\). What is \emph{not} in that paper,
and should be emphasized, is the finite-time dump calculus and the
\(I\)-economy.

\paragraph{Emphasize.}
\begin{enumerate}
\item Dump process \(W_t\) (R17, R18, R20, R22). KT give the
  \emph{eventual} killing position \(K=X(T_H)\). The chapter keeps a
  time-dependent released count \(W_t=K\mathbf{1}_{\{t\ge\tau\}}\) and
  closes \(V(t)\), \(\mathbb{E}[W_t^2]\), \(\mathrm{Var}(W_t)\), and
  current yield as functions of \(I(t)\). These were not found. They
  matter because Chapter~6's age kernels are \(\delta K(\alpha)\) and
  \(V(\alpha)\).
\item Load moments as polynomials in \(I\) (R23). Equivalent moments live
  in the PGF. The named proposition (degree \(p+1\), roots at \(a\) and
  \(b\), leading coefficient \((-1)^p p!(\lambda/\delta)^p\)) was not
  found. Every later kernel is then a polynomial in one monotone clock.
\end{enumerate}

\paragraph{Cite as known (KT 1982).}
Process with cemetery \(H\) (generator (0.4), §2); roots \(a,b\) (2.2);
linear-fractional PGF (2.5); \(I(t)=P(T_H>t)\) (2.7); geometric killing
position, mean \(a/(a-1)\) (2.8); geometric QSD (2.9), uniqueness by
Zhang--Zhu 2011. Also already: \(J\) and \(K\) as derivatives of the PGF
(so Chapter~1's ``second moment beyond the literature'' is false).
Brockwell--Gani--Resnick 1982 is the right historical citation but is
\emph{more general}.

\itemtable
\itemhead
R01 & Linear BDC with release & \alr &
KT 1982; Brockwell et al.\ 1982 more general \\
R02 & Killing vs catastrophe (\(H\) vs \(0\)) & \alr &
KT already keep \(H\); Baake et al.\ 2024 revisit the naming \\
R03 & Riccati for \(I\) & \alr & KT; linear BD Riccati (Kendall) \\
R04 & Characteristic roots \(a,b\) & \alr & KT (2.2) \\
R05 & Vieta / gap identities & \alr & Elementary from the quadratic \\
R06 & \(I(t)\) closed form & \alr & KT (2.7) \\
R07 & \(b=\) extinction-before-rupture & \alr &
KT \(P(\text{hit }0\text{ before }H)\) \\
R08 & \(D(t)\) internal extinction & \alr & Split of the PGF at \(z=0\) \\
R09 & \(\Ifix=I-D\) & \alr & KT two conditionings \\
R10 & \(\Ifix\) initial slope & \alr & Differentiation of R09 \\
R11 & Multi-founder \(I^k-D^k\) & \alr & Branching property \\
R12 & Non-fixation does not factorise & \alr & Same \\
R13 & \(\dot I=-\delta J\) & \alr & Flux / Green \\
R14 & Mean load \(J\) & \alr & \(\partial_z G(1,t)\) \\
R15 & Open \(J\)--\(K\) hierarchy & \alr & Standard moment ODEs \\
R16 & Second moment \(K\) & \alr &
\(\partial_{zz}G\); Chapter~1 over-claim \\
R17 & Mean release \(V(t)\) & \ext &
Eventual \(V_\infty\) is KT; finite-time \(V(t)\) not found \\
R18 & Release second moment & \ext &
Eventual geometric second moment is KT \\
R19 & \(\mathrm{Var}(X)\) & \alr & \(K-J^2\) \\
R20 & \(\mathrm{Var}(W)\) & \ext & Finite-time dump variance not found \\
R21 & Exclusion / \(\mathrm{Cov}(X,W)=-JV\) & \nar &
Definitional (\(XW\equiv 0\)) \\
R22 & Current yield \(\mathbb{E}[X+W]\) & \ext &
Bookkeeping in the \(W\) calculus \\
R23 & Load moments polynomials in \(I\) & \ext &
Organisation of the PGF; named proposition not found \\
R24 & \(I\) as clock & \nar & Time change of an autonomous ODE \\
R25 & Burst-size first two moments & \alr &
KT geometric, mean \(a/(a-1)\) \\
R26 & Geometric burst preview, \(\mu=0\) & \alr & Yule + killing; KT \\
R27 & Burst-time density \(\delta J\) & \alr & Flux \\
R28 & Productive conditional mean plateau & \alr &
QSD mean (KT / Zhang--Zhu) \\
R29 & \(J/\Ifix\) vs \(J/I\) & \alr & Two conditionings \\
R30 & No-death limit & \alr & \\
R31 & No-catastrophe \(\to\) BD & \alr & Kendall \\
R32 & Critical births = deaths & \alr & \\
R33 & No-births limit & \alr & \\
R34 & Four-ODE reconstruction & \nar & Internal check \\
R35 & Landscape of \(b\) and \(V_\infty\) & \nar & Figure \\
R36 & Mean load mixes fates & \nar & Reading of \(J\) \\
R37 & Load SD exceeds mean & \nar & Reading of R19 \\
R38 & Endpoints & \nar & Consistency \\
R39 & Special case of Brockwell & \alr & Correct placement \\
R40 & \(I_\pm\) notation & \nar & Dictionary \\
R41 & Non-explosion & \alr & Linear rates; Cairns--Pollett Thm~1 \\
R42 & Three research questions & \nar & Framing \\
\end{longtable}

\paragraph{Suggested wording.}
\begin{quote}
The linear killing process of Karlin and Tavar\'e (1982) is the
intracellular model. Their generating function already supplies \(I\),
\(J\), \(K\), and the eventual burst-size moments. The present chapter
adds the finite-time released count \(W_t\) and shows that every load
moment is a polynomial in \(I\). Those two additions are what the
population chapter consumes.
\end{quote}

\paragraph{Why it matters biologically.}
A lytic cell (\yp{} in a macrophage, \ft{}, lytic phage) does not drip.
Per-capita catastrophe \(\delta n\) makes a crowded cell more likely to
burst. Distinguishing internal clearance (\(D\)) from rupture (\(1-I\))
is the difference between an empty intact cell and a dump. \(V(t)\) is
the expected dump by age \(t\); that is the object a within-host model
needs, not a fitted constant \(p\).

% =====================================================================
\chapter{Chapter 5 --- BDC distribution, QSD, burst statistics}
\label{ch:5}

\paragraph{Source.} \texttt{CH5\_REWRITE/sections/}. 49 items.

\paragraph{Scope.}
Complete law of one bursting cell: generating function, sliding geometric
occupancy, QSD, burst-size identity, burst time, conditional means,
several founders, chained transfer, single-cell bursting vs budding. The
chapter organises itself around one identity: the burst-size law equals
the quasi-stationary law.

\paragraph{Verdict.}
That identity is Karlin and Tavar\'e (1982), eqs.~(2.8)--(2.9). Zhang and
Zhu (2011) restate uniqueness of the geometric QSD. Kendall (1948)
already has the unkilled sliding geometric. Appendix rows marked N for
\(P(t)\), \(p_n\), QSD, burst law, and \(\mathbb{E}[T_{\mathrm{prod}}]\)
should be S (standard) or R (recap of Chapter~4). What to emphasize is
the proof, the timing laws, the multi-founder yield, and the
chained-transfer laboratory.

\paragraph{Emphasize.}
\begin{enumerate}
\item Telescoping / size-bias criterion (R22). New \emph{proof} that
  burst = QSD, and a characterisation (linear hazard + sliding geometric
  + \(p_1=P'/\lambda\)).
\item \(\mu=0\) harmonic / Gumbel timing (R26--R27). \(\tau\mid\mathcal{K}=n\)
  as a sum of exponentials; Gumbel limit. Budding comparator is linear in
  \(n\).
\item Multi-founder kernels (R31--R34). \(g_k\); free-sum trap;
  \(\mu=0\) lifetime yield \(k+\lambda/\delta\) (sub-extensive
  unconditionally). Several founders sharing one cell do not add solo
  yields.
\item Chained immediate transfer (R35--R42). Wholesale transfer of the
  load. Negative-binomial rupture sizes are textbook sums of geometrics;
  the laboratory (and fragility as soon as \(\mu>0\)) was not found under
  this name.
\item Three conditional burst means / late-burst factor (R23--R24, R29).
  Size-bias of the geometric QSD. ``Factor of two'' is \(a\to 1\), not
  universal.
\item Single-cell BDC vs budding (R45--R46) and \ft{}/anthrax rows (R48).
  Pearson 2011 already has mean-matched burst vs continuous at population
  level with fixed \(N\). Numerical timings at Carruthers \ft{} and
  Williams anthrax rates are new evaluations.
\end{enumerate}

\paragraph{Do not claim as new.}
Sliding geometric occupancy (R08, R11, R12) --- KT (2.5). QSD geometric
(R14) --- KT (2.9). Burst size = QSD (R20, R21) --- KT (2.8), the
headline. Mean productive lifetime (R18) --- KT Green function
\(\lambda^{-1}\log(a/(a-1))\). Recap of Chapter~4 closed forms.

\itemtable
\itemhead
R01 & BDC process & \alr & KT 1982 \\
R02 & Joint process / \(W_t\) & \alr & KT killing position; Ch.~4 \\
R03 & Roots, AB identity & \alr & Ch.~4 / KT \\
R04 & \(I,D,\Ifix\) & \alr & Ch.~4 / KT \\
R05 & \(J,K,V\), conditional mean burst & \alr & Ch.~4 / KT \\
R06 & Second moment of one-cell release & \alr &
Ch.~4 (eventual); finite-time is CH4-R18 \\
R07 & Kolmogorov equations, PGF PDEs & \alr & \\
R08 & Single-founder GF & \alr & KT (2.5) linear-fractional \\
R09 & \(k\)-founder GF & \alr & Branching property \\
R10 & Two rupture conventions & \ext & CH4-R02; bookkeeping \\
R11 & State probabilities & \alr & Coefficients of (2.5) \\
R12 & Load is geometric at every time & \alr & KT; Kendall unkilled \\
R13 & Positive states carry \(\Ifix\) & \alr & \\
R14 & QSD & \alr & KT (2.9); Zhang--Zhu 2011 \\
R15 & QSD moments & \alr & Geometric moments \\
R16 & QSD decay rate & \alr & \\
R17 & Decay rate = BDC time constant & \alr & \\
R18 & Mean productive lifetime & \alr & KT Green \\
R19 & Occupation integral of \(p_k\) & \alr & Green function \\
R20 & Burst-size distribution & \alr & KT (2.8) \\
R21 & Burst size = QSD & \alr & \textbf{KT identity} \\
R22 & Telescoping criterion & \ext & New proof / characterisation of R21 \\
R23 & Size-biased burst given rupture time & \ext & Geometric size-bias \\
R24 & Late-burst mean & \ext & \\
R25 & Defective burst-time density & \alr & Flux \(\delta J\) \\
R26 & \(\tau\mid k\), \(\mu=0\) & \ext & Yule / exponential order stats \\
R27 & Conditional CDF, Gumbel limit & \ext & \\
R28 & Mean burst time on bursting & \ext & Evaluation \\
R29 & Three conditional burst means & \ext & Circularity warning \\
R30 & \(k\)-founder fixation & \alr & Branching property \\
R31 & \(k\)-founder moments, \(g_k\) & \ext & \\
R32 & Superlinear flux, free-sum trap & \ext & \\
R33 & Lifetime yield from \(k\) founders & \ext & \\
R34 & \(\mu=0\) yield, sub-extensivity & \ext & \(k+\lambda/\delta\) \\
R35 & Chained immediate-transfer model & \ext & Not found as this laboratory \\
R36 & Event types ignore the load & \alr & Competing exponentials \\
R37 & NB rupture-size law & \ext & Textbook NB applied here \\
R38 & Chained rupture-size moments / PGF & \ext & \\
R39 & Chained rupture times & \ext & \\
R40 & Inter-rupture Laplace & \ext & Phase-type \\
R41 & Mean inter-rupture interval & \ext & \\
R42 & Second rupture-time Laplace & \ext & \\
R43 & Coefficient extraction, general \(\mu\) & \ext &
Ch.~2 method applied; \(\mu>0\) open \\
R44 & Simulation verification & \nar & \\
R45 & Bursting vs budding structural table & \ext &
Pearson/Yuan population, fixed \(N\) \\
R46 & Matching means \(\neq\) matching support & \ext &
Pearson's logical point, single-cell arithmetic here \\
R47 & Five limiting-dilution predictions & \ext &
Translation of R14/R21 into an assay \\
R48 & \ft{} / anthrax numerical extremes & \ext &
Carruthers 2020; Williams 2021 rates \\
R49 & Geometric robust to \(k\), fragile to chaining with death & \ext &
Synthesis \\
\end{longtable}

\paragraph{Suggested wording.}
\begin{quote}
Karlin and Tavar\'e (1982) already proved that for linear birth--death
with killing, occupancy is linear-fractional, the QSD is geometric, and
the killing position is that same geometric. The present chapter gives a
size-bias proof of that identity, the burst-time law at \(\mu=0\), the
multi-founder yield \(k+\lambda/\delta\), and a chained-transfer
laboratory in which the geometric structure fails as soon as death is
present. Those are the distributional facts a lytic-cell assay can test.
\end{quote}

\paragraph{Why it matters biologically.}
A fitted mean burst size does not determine establishment (Pearson). For
a genuinely lytic pathogen the whole geometric family is determined by
\((\lambda,\mu,\delta)\), and the same three rates span \ft{} (rare
clearance, huge bursts, low ID) and anthrax (almost-sure clearance, tiny
bursts, high ID). Several co-infecting founders sharing one cell do not
add solo yields.

% =====================================================================
\chapter{Chapter 6 --- Burst-aware renewal and population comparison}
\label{ch:6}

\paragraph{Source.} \texttt{CH6\_REWRITE/sections/}. 33 items.

\paragraph{Scope.}
What replaces BMVR's \(pI\) when release is terminal and batched.
Burst-aware renewal forced by BDC kernels; classical BMVR as an
exponential-phase projection; identifiability; \(R_0\); flooding
parameter \(L\); growth--extinction tradeoff.

\paragraph{Verdict.}
This chapter contains two of the four original statements in the thesis
(the \(L\)-flooding identity and the \(L=1\) coincidence of laws), sitting
on top of a renewal system that is Nelson et al.\ (2004) and an \(R_0\)
invariance that is Pearson (2011), first paragraph. Write the skeleton as
an application of Nelson; write \(R_0\) as Pearson; lead with \(L\).

\paragraph{Emphasize.}
\begin{enumerate}
\item Flooding identity (R20) --- original as a statement.
\[
z_{\mathrm{ext}}^{\mathrm{burst}}-z_{\mathrm{ext}}^{\mathrm{bud}}
=(L-1)\Bigl(\frac{1}{m}-1\Bigr),\qquad L:=a(1-b).
\]
Bursting lowers establishment extinction iff \(L>1\), i.e.\
\(1/\delta<1/\lambda+1/\mu\) (\(\mu=0\) always). Pearson 2011 cannot flip
this sign; Quirouette et al.\ 2023 already said extinction tracks the
offspring \emph{law}, not timing; Komarova 2007 floods antibodies
(different object).
\item \(L=1\) coincidence of laws (R21). When \(L=1\) the two offspring
  laws are the same geometric. Variance order:
  \(\mathrm{Var}_{\mathrm{bud}}-\mathrm{Var}_{\mathrm{burst}}\propto(L-1)\).
\item Kernels from BDC, not posited (R04--R06) --- extends Nelson.
  \(S\) and \(g=\delta K\) are closed forms in \((\lambda,\mu,\delta)\).
  No free age schedule to fit. \(p\) is no longer a parameter of the
  organism.
\item Exponential-phase projection (R09--R10) --- extends Euler--Lotka.
  Fitted \(p\) is a \emph{phase} property. Do not claim the projection
  idea is new.
\item Identifiability 3-to-2 (R16) --- extends Nelson/Miao.
\item Growth--extinction tradeoff (R22--R24) --- numerical. General
  \(r_{\mathrm{bud}}>r_{\mathrm{burst}}\) is unproved --- leave it
  unproved.
\end{enumerate}

\paragraph{Do not claim as new.}
BMVR (R01). Exponential kernels \(\equiv\) BMVR (R07). \textbf{\(R_0\)
invariance at matched yield (R14). Pearson 2011 opening paragraph.} Use
it, cite it, do not number it as a new theorem.

\itemtable
\itemhead
R01 & Classical BMVR / budding cell & \alr &
Perelson 1996; Nowak 1996; McLean 1993 \\
R02 & Synchronous cohort is scheduled & \ext & Obstruction to constant \(p\) \\
R03 & Obstruction to constant \(p\) & \ext & \\
R04 & Survival and release kernels & \ext &
Nelson 2004 posited \(P(a)\); here closed-form BDC \\
R05 & Burst-aware renewal system & \ext & Nelson eqs.~(1) \\
R06 & Universal skeleton, any kernels & \ext & Same \\
R07 & Exponential kernels \(\to\) BMVR & \alr & \\
R08 & Renewal verification suite & \nar & \\
R09 & Exponential-phase projection & \ext &
Euler--Lotka / Diekmann / Nelson \\
R10 & \(p_{\mathrm{eff}}\) map endpoints & \ext & \\
R11 & \(p_{\mathrm{eff}}\) decreases in \(r\) (numerical) & \ext & \\
R12 & Generation kernel, \(R_0\), Euler--Lotka & \ext & Diekmann \\
R13 & Overlays: what the projection loses & \ext & \\
R14 & \(R_0\) unchanged by bursting & \alr & \textbf{Pearson 2011} \\
R15 & Continuum \(R_0\) vs discrete \(m\) & \alr & \\
R16 & Identifiability obstruction & \ext & Nelson 2004; Miao 2011 \\
R17 & Fitting-practice predictions & \ext & \\
R18 & Bursting offspring PGF / moments & \ext &
Pearson random-burst PGF; KT size-at-killing \\
R19 & Bursting extinction probability & \ext &
Pearson root of thinned PGF \\
R20 & Flooding identity / \(L\) criterion & \orig &
Pearson (no \(L\)); Quirouette; Komarova different object \\
R21 & \(L=1\) coincidence of laws & \orig & Linear-fractional family \\
R22 & Budding invades faster at matched \(R_0\) & \ext &
Pearson 1.75\,d vs 2.46\,d; Komarova \\
R23 & Generation-time mechanism & \ext & Wallinga--Lipsitch \\
R24 & Growth--extinction Pareto in \(\delta\) & \ext &
Numerical; general order unproved \\
R25 & Five-ODE ladder & \ext & Reach of the construction \\
R26 & BMVR as fast-export / QSS & \ext & \\
R27 & Reset-catastrophe intermediate & \ext & Catalogue \\
R28 & Partial-release mean field & \ext & \\
R29 & HIV-stage renewal and Allee & \ext & Application \\
R30 & HIV mean-linear, not exponential BDC & \alr & Hataye / HIV production \\
R31 & What carries beyond BDC & \nar & Discussion \\
R32 & Three regimes, flooding-boundary numerics & \nar & Figures for R20 \\
R33 & Hypergeometric Laplace of kernels & \ext &
Ch.~4 \(I\)-polynomials in Laplace space \\
\end{longtable}

\paragraph{Suggested wording.}
\begin{quote}
Nelson et al.\ (2004) already replaced constant \(p\) by an age-dependent
production schedule. Here that schedule is not fitted: it is the BDC flux
\(\delta K(\alpha)\). Pearson et al.\ (2011) already showed that matched
lifetime yield equalises \(R_0\) while leaving stochastic extinction
free. Substituting the BDC burst law yields
\(\Delta z_{\mathrm{ext}}=(L-1)(1/m-1)\). When \(L=1\) the two offspring
laws coincide. That is the chapter's original comparison; the renewal
skeleton and the \(R_0\) lemma are inherited.
\end{quote}

\paragraph{Why it matters biologically.}
A fitted HIV \(p\) is not a property of \yp{}. For a lytic pathogen,
release and death are the same event, the extracellular curve is a
convolution against a closed-form kernel, and whether bursting helps
establishment is a criterion on \((\lambda,\mu,\delta)\), not a slogan.
The same criterion will reverse in Chapter~8 once the immune functional
is a threshold rather than linear thinning.

% =====================================================================
\chapter{Chapter 7 --- Two-type BDC}
\label{ch:7}

\paragraph{Source.} \texttt{CH7/Two\_Type\_Chapter/document MAIN/sections/}.
31 items.

\paragraph{Scope.}
Triangular two-type birth--death (type 1 \(\to\) type 2 conversion) plus
a global absorbing catastrophe at rate \(\delta_1 X+\delta_2 Y\).
Observables: no-catastrophe probabilities \(S(t)\) and \(G(t)\), not
terminal counts. Type 2 is autonomous; type 1 is a forced Riccati that
linearises to Gauss hypergeometric. Biological reading: early vs adapted
intracellular \yp{}.

\paragraph{Verdict.}
This chapter has the other two original statements (exact \(S\), and the
killed PGF that recovers Antal--Krapivsky). The \emph{method} is Antal and
Krapivsky (2011). Type-2 \(G\) is one-type linear BDC (Karlin--Tavar\'e /
Chapter~4). Write the method as inherited; write the object as original.

Verified from AK arXiv:1105.1157: their backward system is this chapter's
ODEs at \(\delta_1=\delta_2=0\) with PGF data \(x,y\). Survival data
\(S(0)=G(0)=1\) then force \(S\equiv G\equiv 1\). So \(\delta=0\) in
\(S,G\) does \emph{not} recover AK. Recovery holds for the killed PGF.
Li (2024) ``two-type catastrophes'' is a homonym. No published
\({}_2F_1\) form for type-specific linear catastrophe \(\delta_1 x+\delta_2 y\)
was found.

\paragraph{Emphasize.}
\begin{enumerate}
\item Exact \(S(t)\) (R06) --- original as a statement. Wronskian mix of
  a real-axis \({}_2F_1\) pair along \(z_2<0\),
\[
\widehat S(\hat t)=h-\hat q_2 z_2(\hat t)
\frac{K_U U'(z_2)+K_V V'(z_2)}{K_U U(z_2)+K_V V(z_2)},
\]
with \(U={}_2F_1\), \(V=(-z)^{1-C}{}_2F_1\). Nearest miss: AK
terminal-count PGF. A lineage that grew early and died can contribute
more catastrophe exposure than a small lineage that remains alive.
\item Killed PGF (R23) --- original as a companion. Same mix, data
  \(x,y\). AK is \(\delta_i=0\) of \emph{this} object.
\item Method (R05) --- extends AK. Cite AK 2011 and Bender--Orszag. Do
  not claim a new special-function theory.
\item \(\lambda_2=0\) Bessel boundary (R18) --- extends. Ordinary
  \(J_\rho,Y_\rho\) with \(\delta_2>0\) is not AK's modified-Bessel
  bi-critical PGF.
\end{enumerate}

\paragraph{Do not claim.}
Type-2 \(G\), \(\delta_2=0\) type-1, \(\nu=0\) (R04, R15, R19): one-type
linear BDC. Under-cited. Point at Karlin--Tavar\'e / Chapter~4 \(I(t)\).
The biological EQ/MAT/GATE/EARLY comparison as a theorem (MAT is a
working hypothesis).

\itemtable
\itemhead
R01 & Two-type BDC process & \ext & AK skeleton + KT/Brockwell killing \\
R02 & Backward equations & \alr & Standard first-step / branching \\
R03 & Autonomous type-2 & \ext & Triangular structure is AK \\
R04 & Closed form for \(G\) & \alr & KT / Ch.~4 \(I(t)\) (one-type BDC) \\
R05 & Gauss hypergeometric reduction & \ext & AK §§II--III \\
R06 & Exact finite-time \(S\) & \orig & AK PGF is a different object \\
R07--R11 & Derivation steps / branch management & \nar & Method of R06 \\
R12 & Forever-avoidance \(S\to h\) & \ext & \\
R13 & \(G\to r_{2,-}\) & \ext & \\
R14 & Equal-rate gap (no independent form) & \ext & Fair negative remark \\
R15 & \(\delta_2=0\) elementary Riccati & \alr & One-type \\
R16 & \(\lambda_1=0\) linear integral & \alr & \\
R17 & Limits catalogue & \ext & \\
R18 & \(\lambda_2=0\) Bessel mix & \ext & Not AK's \(I_\nu,K_\nu\) \\
R19 & \(\nu=0\) decoupled & \alr & Two independent one-type BDCs \\
R20--R21 & Parameter dictionary / scaling & \alr & \\
R22 & Killed-PGF backward system & \ext & AK (10)--(11) at \(\delta=0\) \\
R23 & Killed-PGF closed form & \orig & Recovers AK at \(\delta_i=0\) \\
R24 & Ultimate counts on catastrophe-free paths & \ext & \\
R25 & Ultimate release law & \ext &
KT size-at-killing, two-type packaging \\
R26 & Interpretation stub & \nar & Empty file \\
R27 & Biological application (four orderings) & \ext &
EQ/MAT/GATE/EARLY: reading, not identification \\
R28--R29 & Companion biology extra arguments & \nar & \\
R30 & Early vs adapted YP phenotypes & \alr & Ke 2013; Pujol--Bliska \\
R31 & MAT bursting as working hypothesis & \alr &
Monack 1997; Peters 2013; not a theorem \\
\end{longtable}

\paragraph{Suggested wording.}
\begin{quote}
Antal and Krapivsky (2011) solved the triangular two-type birth--death
process for terminal population counts by reducing a Riccati system to
Gauss's equation. The present chapter keeps that triangular skeleton and
that reduction, and changes the observable: a type-specific linear
catastrophe clock, and the probability that it has not yet fired. Setting
the catastrophe rates to zero in those survival probabilities does not
recover Antal--Krapivsky; it recovers them from the killed generating
function, which is the same mix with PGF data. Type 2 is the classical
one-type killing probability.
\end{quote}

\paragraph{Why it matters biologically.}
Intracellular \yp{} is not one type. Conversion (adaptation, Yop
expression) and loss of macrophage containment can run at different
per-capita rates. \(S(t)\) is the chance that a lineage started as the
early type has not yet blown the cell, integrating exposure along the
whole path --- the right object if catastrophe is ``being seen'' rather
than ``being extinct.''

% =====================================================================
\chapter{Chapter 8 --- Bursting and budding pathogenesis}
\label{ch:8}

\paragraph{Source.} \texttt{CH8\_REWRITE/sections/}. 21 items.

\paragraph{Scope.}
Does bursting help once immune capacity is depletable? Schedule-alone
cannot matter (linear rates). Two nonlinear models: intracellular
replicator--suppressor; extracellular fixed removal. Closed forms, a
certainty threshold, \(\mathcal{R}(\vartheta)=V_\infty a^{-\vartheta}\),
and reversal of Chapter~6's flooding criterion. The chapter correctly
says that comparing bursting and budding is not new (Pearson, Komarova).

\paragraph{Verdict.}
No original statements in the strict sense after adversarial checking:
every New mark has a mechanism neighbour (Pilyugin--Antia, geometric
stop-loss, Yuan--Allen ranking-switch, Komarova flooding). The
constructions and the \(L\)-reversal under a threshold functional are
still the chapter's contributions. Cite Yuan 2011 and Pilyugin--Antia
2000, which the chapter currently misses.

\paragraph{Emphasize.}
\begin{enumerate}
\item Reversal of the flooding criterion (R17) --- strongest extension.
  Same matched-mean pair as Chapter~6; functional
  \(x\mapsto(x-\vartheta)_+\) is convex, so the Chapter~6 variance order
  \(\propto(L-1)\) \emph{flips} the sign: bursting now helps when
  \(L<1\). Yuan and Allen (2011) already reverse budding vs bursting
  under an HIV-style mass-action immune term. Different mechanism,
  different criterion, but the qualitative published fact is the
  ranking-switch. The thesis contribution is the complementary \(L\)
  criterion on the \emph{same pair of laws} as Chapter~6.
\item \(\mathcal{R}(\vartheta)=V_\infty a^{-\vartheta}\) and
  \(\vartheta^\ast=\log V_\infty/\log a\) (R11--R12). Geometric stop-loss
  is actuarial folklore. Not found as a BDC burst + integer immune
  budget.
\item Replicator--suppressor CTMC (R03) and certainty diagonal
  (R05--R06). Pilyugin--Antia 2000: finite macrophage pool that
  \emph{recycles}; their \(P^\ast\) depends on rates. The
  consume-not-recycle choice makes \(i>j\Rightarrow\) leftover
  replicators even with no births --- combinatorial, rate-independent.
\item Two flooding arguments (R18). Batch-count is Komarova's intuition
  and is elementary. Distributional (two laws, one toll) is R17.
\end{enumerate}

\paragraph{Do not claim.}
Bursting vs budding as a question. Finite killing budget as a new
mechanism. Schedule-alone invariance (R01) as new. YP two-phase as
\(\vartheta^\ast\)-crossing (R20): Ke 2013 answers the switch with
temperature / pH / YCV. Replace the flooding placeholder cite with
Komarova (2007).

\itemtable
\itemhead
R01 & Schedule-alone cannot matter & \alr &
Pearson mean-field; Ch.~6 \(R_0\); Diekmann \\
R02 & Chained-rupture process & \alr & Ch.~5 laboratory; successive BDC \\
R03 & Replicator--suppressor process & \ext &
Pilyugin--Antia 2000 (recycle); Lanchester \\
R04 & Replicator--suppressor numerics & \ext & \\
R05 & Certainty threshold \(i>j\) & \ext &
Combinatorial; Pilyugin vulnerability; Li 2002 neutrophil threshold \\
R06 & Diagonal law \(\pi_{j,j}\) & \ext &
Pochhammer product; \(1/(j+1)\) at \(\nu=1\) classical \\
R07 & Off-diagonal / simulations & \ext & \(\pi_{i,j}\) open in closed form \\
R08 & Detection avoidance (named, not modelled) & \nar &
Conway--Ribeiro 2018 \\
R09 & Fixed-removal model & \ext &
Komarova flooding; complement consumption; integer-\(\vartheta\)
caricature \\
R10 & Geometric burst law (inherited) & \alr & KT / Ch.~4 / Ch.~5 \\
R11 & \(\mathcal{R}(\vartheta)=V_\infty a^{-\vartheta}\) & \ext &
Geometric stop-loss \\
R12 & \(\vartheta^\ast=\log V_\infty/\log a\) & \ext & Corollary of R11 \\
R13 & Subcritical regime created by \(\vartheta\) & \alr &
\(V_\infty>1\) inherited; commentary \\
R14 & Numerical \(\vartheta^\ast\) rows & \ext & Illustration \\
R15 & Subductive proliferation (named, not modelled) & \nar & \\
R16 & Ch.~6 \(R_0\) / \(L\) recalled & \alr & Ch.~6 \\
R17 & Flooding criterion \emph{reverses} & \ext &
Yuan 2011 ranking-switch (different immune); convex order of Ch.~6 \\
R18 & Two flooding arguments & \ext & Komarova = batch-count \\
R19 & Selection on the bursting--budding spectrum & \ext &
Komarova; Wang lysis-timing; Ch.~6 spectrum \\
R20 & YP two-phase as threshold-crossing & \ext\ (weak) &
Ke 2013 temperature cue, not \(\vartheta^\ast\) \\
R21 & Verification tables & \nar & 0 sign errors on reversal \\
\end{longtable}

\paragraph{Suggested wording.}
\begin{quote}
Pearson, Komarova, and Yuan and Allen already compared bursting and
budding. What is missing from those comparisons, and from Chapter~6's
linear-clearance flooding theorem, is a finite killing budget. Two
caricatures are introduced. Intracellularly, unreplenished one-to-one
annihilation (cf.\ Pilyugin and Antia, who recycle killers) produces a
rate-independent certainty threshold. Extracellularly, an integer toll
on the geometric BDC burst yields
\(\mathcal{R}(\vartheta)=V_\infty a^{-\vartheta}\). Because
\(x\mapsto(x-\vartheta)_+\) is convex, Chapter~6's \(L\) criterion
reverses. Yuan and Allen (2011) already found a ranking-switch under
mass-action immunity; the present reversal is the complementary
criterion on the same BDC/budding pair.
\end{quote}

\paragraph{Why it matters biologically.}
Macrophage granules and extracellular complement are depletable. Linear
clearance (Chapter~6) cannot see a safety-in-numbers threshold. Whether
bursting helps then depends on \emph{which functional} one puts on the
same two offspring laws --- the point of R17--R18, and the reason
Komarova's unconditional flooding must not be copied.

% =====================================================================
\chapter{Chapter 9 --- Variability in evolution}
\label{ch:9}

\paragraph{Source.} \texttt{CH9/document MAIN/sections/}. 17 items.
No numbered theorem environments.

\paragraph{Scope.}
High early diversification after ecological relief as a mechanistic
variable-rate phenomenon, not only Nee's push-of-the-past sampling
artefact. Core: no-death time-inhomogeneous Yule whose speciation rate
tracks logistic \emph{individual} abundance. Then speculative models:
Pimentel genetic feedback, rise--run--ruin--rejuvenation, birth--plague,
vitality.

\paragraph{Verdict.}
No original statements. Diversity-dependent diversification (DDD) and
time-dependent Yule/birth--death are a mature literature. The chapter's
claim that ``nobody has published a logistic version'' is false
(Sepkoski; Rabosky DDL; Etienne--Haegeman DDD). A no-death Yule cannot
test POTPa, which requires extinction and conditioning on survival. Pull
of the present is title-only.

\paragraph{Emphasize (weakly).}
\begin{enumerate}
\item Logistic-individual Yule (R03). \(N(t)\) is a deterministic
  individual logistic; \(Y\) is an ordinary inhomogeneous Yule (Kendall
  1948). DDD puts \(\lambda(n)\propto 1-n/K\) on \emph{species} count.
  The two-layer unused-capacity coupling is the only residual
  distinction. Not a new process class.
\item Pimentel Moran + minority-resistance kernel (R08--R10). Still in
  the Pimentel--Levin--Pease genetic-feedback class. Simulations show
  reversal-to-fixation, not Pimentel's missing oscillation. Unfinished.
\item RRRR and birth--plague (R13, R15). Assemblies of BDC+immigration
  and SI-with-sparking. Present as exploratory models, not results.
\end{enumerate}

\paragraph{Do not claim.}
A new alternative to Nee 1994 (the last two decades of DDD vs
time-dependent BD vs POTPa \emph{are} that alternative). The explicit
mean \(e^{\sigma Kt}((\xi+1)/(\xi+e^{rt}))^{\sigma K/r}\) as a named
discovery (three-line specialisation of Kendall). PGF / geometric
occupancy (R06--R07): Kendall. Punctuated-equilibria theorem (R01).

\itemtable
\itemhead
R01 & Relaxed-selection punctuated equilibria & \ext\ (hypothesis) &
Eldredge--Gould; Newman--Cohen--Kipnis 1985; Yoder; Lahti \\
R02 & POTPa vs variable-rate origination & \alr\ (as a question) &
Nee et al.\ 1994; Budd--Mann 2018; DDD literature \\
R03 & Logistic-speciation Yule & \ext\ (barely) &
Kendall TD-Yule; Sepkoski; Rabosky/Etienne DDD on species; Hubbell \\
R04 & Ensemble slowdown figure & \alr\ (as a pattern) &
Generic to declining \(\lambda\); Pannetier TD\(\equiv\)DD for LTT \\
R05 & Logistic-Yule mean & \alr &
Kendall \(\exp\int\lambda\); elementary integral \\
R06 & Logistic-Yule PGF & \alr & Kendall inhomogeneous Yule PGF \\
R07 & Occupancy distribution & \alr & Geometric / Kendall \\
R08 & Pimentel Moran+resistance & \ext & Pimentel; Levin; Pease \\
R09 & Reversal and divergence & \ext & Reversal shown, not oscillation \\
R10 & Critical saddle & \ext & \\
R11 & Cellular automaton & \ext & Illustration of R08 \\
R12 & Pimentel + variation collapse & \ext & Unfinished \\
R13 & Rise--run--ruin--rejuvenation & \ext & Assembly of known pieces \\
R14 & RRRR cyclic realisations & \nar & Figures \\
R15 & Birth--plague model & \ext & BDC+immigration; SI with demography \\
R16 & Birth--plague realisations & \nar & Figures \\
R17 & Vitality birth--plague & \nar & Vitality never enters the rates \\
\end{longtable}

\paragraph{Suggested wording.}
\begin{quote}
Nee, May and Harvey (1994) showed that apparent early bursts can be a
sampling artefact of constant-rate birth--death conditioned on survival.
Diversity-dependent and time-dependent speciation models are the existing
mechanistic alternatives. The model used here is a no-death Yule whose
rate tracks unused \emph{individual} capacity (Kendall 1948 with a
logistic clock). It cannot adjudicate push of the past, which requires
extinction. It is offered as a simple variable-rate illustration, not as
a new alternative to that literature.
\end{quote}

\paragraph{Why it matters biologically.}
The thesis's interest is whether \yp{} (and large clades generally) can
show genuinely elevated origination after ecological relief, rather than
only a conditioned-BD illusion. That question is real (Budd and Mann).
This chapter does not yet have a model that can answer it, because the
analytic core has no death.

% =====================================================================
\chapter{Chapter 10 --- Conclusion: evolved strategy of \yp{}}
\label{ch:10}

\paragraph{Source.} \texttt{CH10/document MAIN/sections/}. 15 items.
No numbered theorems. (Source title: Chapter~8.)

\paragraph{Scope.}
Why \yp{} uses a two-phase invasion: macrophage then extracellular /
anti-phagocytic. Conserved-replicative-density mean invariance as a foil;
two-phase within-host branching embedded in a host-to-host
Galton--Watson; ``quadrant'' death-rate argument and a
neutrophil-fraction threshold; speculative public-good accumulation
inside a macrophage.

\paragraph{Verdict.}
No original statements. The biology of the switch is Cavanaugh /
Pujol--Bliska / Ke 2013 (temperature, phagosomal pH, YCV --- not a
neutrophil fraction). The nested BDC\(\to\)BD\(\to\)GW chain is an
application of Antia--Pearson--Gilchrist nested epidemiology plus
Kendall. The public-good model is flagged speculation. Emphasize as
synthesis and application. Do not present \(p_N\) as an evolutionary
theorem.

\paragraph{Emphasize (as application).}
\begin{enumerate}
\item Macrophage boost is not conserved-potential (R03). Appendix-B / R01
  says front-loading birth at fixed total potential does not raise
  expected size. The plague strategy is \emph{not} that comparison: the
  macrophage is a different environment. R01 does not kill RQ1.
\item Nested BDC \(\to\) slightly supercritical BD \(\to\) binary
  host-to-host GW (R06, R08). Architecture is Antia et al.\ 2003 /
  Pearson / Gilchrist. Algebra is Kendall \((\nu/\beta)^N\) mixed over
  the burst size. The displayed \(\mathbb{E}[Z]\) likely has a
  parenthesis error (should be \(2(1-p_0)\)).
\item Public-good accumulation (R13, R15). Neighbours: Asker et al.\
  2023; siderophores / yersiniabactin. Known resistance is PhoP/Q, RovA,
  YCV. Inoculum \(K\) (R14) is not computed. Keep as a possible
  mechanism, not a result.
\end{enumerate}

\paragraph{Do not claim.}
Mean invariance under conserved replicative density (R01) --- Kendall;
Pearson conserved-yield. Two-phase pathogenesis (R04, R11) --- Cavanaugh
1959; Ke 2013. Neutrophil-fraction threshold \(p_N\) (R12): elementary
rearrangement; Ke already answers the switch with temperature; reduced
form as written has a \textbf{sign error} under
\(\mu_N>\mu_E>\mu_M\). Intended threshold is
\((\mu_E-\mu_M)/(\mu_N-\mu_M)\). Two-phase PGF (R07) --- branching
property.

\itemtable
\itemhead
R01 & Conserved-replicative-density mean invariance & \alr &
Kendall 1948; Pearson 2011 conserved yield \\
R02 & Exceptions (non-conservation; Matthew effect) & \nar &
Named, not developed \\
R03 & Macrophage boost \(\neq\) conserved-potential & \ext\ (remark) &
Cavanaugh; Pujol--Bliska; Ke 2013 \\
R04 & Two-phase pathogenesis argument & \alr & Same \\
R05 & Phage-adaptation decaying extracellular death & \ext &
Speculative mechanism \\
R06 & Host-to-host GW embedding & \ext\ (application) &
Whittle/Bartlett; Antia 2003; Pearson; Gilchrist--Sasaki \\
R07 & Two-phase PGF & \alr & Branching property \\
R08 & BDC-to-BD extinction and \(\mathbb{E}[Z]\) & \ext\ (weak) &
Kendall \((\nu/\beta)^N\); check parenthesis \\
R09 & Four within-host variants & \ext & Comparison, not a theorem \\
R10 & Numerical: \(\delta\) controls burst size and \(\mathbb{E}[Z]\) &
\alr\ (as a lesson) & Haldane/Kendall near-critical founder plot \\
R11 & Quadrant rate ordering & \alr & Standard \(\mu_N>\mu_E>\mu_M\) \\
R12 & Neutrophil-fraction threshold \(p_N\) & \alr\ (caricature of Ke) &
Ke 2013 temperature cue; \textbf{sign error} \\
R13 & Public-good accumulation model & \ext\ (speculation) &
Asker 2023; siderophores; local public goods \\
R14 & Critical inoculum \(K\) & \nar & Not computed \\
R15 & Confined public good as group advantage & \ext\ (speculation) &
Nadell/Drescher/Foster \\
\end{longtable}

\paragraph{Suggested wording.}
\begin{quote}
The mean-invariance result of the branching appendix is Kendall's: at
conserved replicative potential, front-loading birth does not raise
expected size. That comparison is the wrong one for \yp{}. The macrophage
is a different environment, not a rearrangement of the same potential
(Cavanaugh 1959; Ke, Chen and Yang 2013). Embedding a BDC intracellular
burst in a slightly supercritical extracellular birth--death process, and
that in a host-to-host Galton--Watson process, is the standard
nested-epidemic construction (Antia et al.\ 2003; Pearson et al.\ 2011)
applied to this pathogen. The neutrophil-fraction threshold is a
caricature of a switch the experimental literature attributes to
temperature and vacuolar cues, not to leukocyte proportions.
\end{quote}

\paragraph{Why it matters biologically.}
RQ1 of the thesis. The strategy is real and famous. The mathematical work
that bears on it is: (i)~near-critical founder sensitivity (Chapters~1
and~10, classical); (ii)~BDC burst size as the macrophage dump
(Chapters~4--5, mostly KT plus \(W_t\)); (iii)~whether bursting then
helps against depletable neutrophils/complement (Chapter~6 \(L\),
Chapter~8 reversal). The conclusion should point at those, not at
\(p_N\).

\end{document}
